\documentclass{article}
\usepackage[utf8]{inputenc}
\usepackage[margin=1in,left=1.5in,includefoot]{geometry}
\usepackage{booktabs}
\usepackage{float}
% Header & Footer Stuff

\usepackage{fancyhdr}
\usepackage{wasysym}
\usepackage{amsmath}
\usepackage{tikz}
\usepackage{tabularx}
\usepackage{amssymb}
\usepackage{graphicx}
\pagestyle{fancy}
\lhead{MATH 220/199}
\rhead{April 14, 2026}
\fancyfoot{}
\lfoot{David Gallardo}
\fancyfoot[R]{}
%027601193
% The Main Document
\begin{document}

\begin{center}
 \LARGE\bfseries MATH 220/199 Week 13 Day 1
 \line(1,0){430}
\end{center}
\section*{Reminders}
\begin{enumerate}
    \item Office hours are held in Davenport 330 at 3pm on Thursdays.
    \item Do the Merit Quiz on the Merit Canvas page.
\end{enumerate}
\section*{What if...Not Rectangle?}
One issue with estimating areas with rectangles is that if a function is very concave/convex, we will get very large overestimates/underestimates. What if we used a new shape, like a trapezoid?
\begin{enumerate}
    \item Let $f(x) = x^2$. Fill in the chart below:
    \newline
    \newline
    $\begin{array}{c|c|c|c|c|c}
        x & 0 & 1 & 2 & 3 & 4 \\
        \hline
        f(x) & & & & & 
    \end{array}$
    \item Compute a Left Riemann sum for the area between $f(x)$ and the x-axis using 4 rectangles.
    \vfill
    \item Compute a Right Riemann sum for the area between $f(x)$ and the x-axis using 4 rectangles.
    \vfill
    \item Now, compute the 4-trapezoidal sum: Construct the trapezoids the same way as the rectangle, but the top side is made by connecting $f(a)$ to $f(a+\Delta)$.
    \vfill
    \item Using $\int\limits_{0}^4 x^2 \text{ d}x = \frac{64}{3}$, which estimate is the best?
    \vfill
\end{enumerate}
\newpage
\section*{Exploring Summation Notation}
The following problems only look hard, they are not actually hard. My goal is to convince you that summation notation is not scary.
\begin{enumerate}
    \item Let $f^{(n)}(x)$ denote the $n^{th}$ derivative of $f(x)$. Say $f^0(x) := f(x)$. Let $f(x) = \sin(x)$. Compute:
    \newline
    $\sum\limits_{i = 0}^8 f^{(2i)}(\frac{\pi}{2}+i\pi)$
    \vfill
    \item Compute $\sum\limits_{n = 0}^{10} (-1)^n \cdot n$
    \vfill
    \item Let $\Delta_{(n, m)}$ be the right triangle with leg lengths $n, m$. Define $\mathbb{A}(\Delta_{(n, m)}) := \text{Area of } \Delta_{(n, m)}$. Compute:
    \newline
    $\sum\limits_{n = 1}^4 \mathbb{A}(\Delta_{(n, n+1)})$
    \vfill
\end{enumerate}

\end{document}
